\subsection{DIRECT SYSTEMS OF MAGMAS}

In this no. and the following we shall assume that I is right directed.

DEFINITION 2. A direct system of magmas relative to the indexing set I is a direct system of sets \((\mathrm{E}_{\alpha}, f_{\beta\alpha})\) relative to I, each \(E_{\alpha}\) having a magma structure and each \(f_{\beta\alpha}\) being a magma homomorphism.

Let \((\mathrm{E}_{\alpha}, f_{\beta\alpha})\) be a direct system of magmas. E will denote the direct limit set \(\lim_{\rightarrow} \mathrm{E}_{\alpha}\) and \(f_{\alpha}\) the canonical mapping of \(\mathrm{E}_{\alpha}\) into E. Recall that

\begin{enumerate}
\def\labelenumi{(\arabic{enumi})}
\setcounter{enumi}{1}
\tightlist
\item
\end{enumerate}

\[
f _ {\beta} \circ f _ {\beta \alpha} = f _ {\alpha} \quad \text { for } \alpha \leqslant \beta ,\tag{3}
\]

\[
\mathrm{E} = \bigcup_ {\alpha \in \mathbf {I}} f _ {\alpha} (\mathrm{E} _ {\alpha}).
\]

By (2), also

\[
f _ {\alpha} (\mathrm{E} _ {\alpha}) \subset f _ {\beta} (\mathrm{E} _ {\beta}) \quad \text { for } \alpha \leqslant \beta .\tag{4}
\]

If \(x_{\alpha}, y_{\alpha} \in \mathrm{E}_{\alpha}\) are such that \(f_{\alpha}(x_{\alpha}) = f_{\alpha}(y_{\alpha})\) , there exists a \(\beta \geqslant \alpha\) such that \(f_{\beta \alpha}(x_{\alpha}) = f_{\beta \alpha}(y_{\alpha})\) .

PROPOSITION 1. There exists on E one and only one magma structure for which the mappings \(f_{\alpha} \colon \mathrm{E}_{\alpha} \to \mathrm{E}\) are homomorphisms. If the magmas \(\mathrm{E}_{\alpha}\) are associative (resp. commutative), so is E. If the magmas \(\mathrm{E}_{\alpha}\) and the homomorphism \(f_{\beta \alpha}\) are unital, so are the magma E and the homomorphisms \(f_{\alpha}\) .

The magmas \(E_{\alpha}\) will be written multiplicatively.

Let \(x, y\) be in E. There exist \(\alpha\) in I and \(x_{\alpha}, y_{\alpha}\) in \(E_{\alpha}\) such that \(x = f_{\alpha}(x_{\alpha})\) and \(y = f_{\alpha}(y_{\alpha})\) . If there exists a magma structure on E for which \(f_{\alpha}\) is a homomorphism, then \(x.y = f_{\alpha}(x_{\alpha}y_{\alpha})\) , whence the uniqueness of this magma structure.

To prove the existence, we must prove that for \(\alpha, \beta\) in I, \(x_{\alpha}, y_{\alpha}\) in E \(_{\alpha}\) and \(x_{\beta}', y_{\beta}'\) in E \(_{\beta}\) , the relations

\[
f _ {\alpha} (x _ {\alpha}) = f _ {\beta} (x _ {\beta} ^ {\prime}), \quad f _ {\alpha} (y _ {\alpha}) = f _ {\beta} (y _ {\beta} ^ {\prime})\tag{5}
\]

imply \(f_{\alpha}(x_{\alpha}y_{\alpha}) = f_{\beta}(x_{\beta}'y_{\beta}')\) . For \(\gamma \geqslant \alpha\) and \(\gamma \geqslant \beta\) , we set \(x_{\gamma} = f_{\gamma \alpha}(x_{\alpha})\) , \(y_{\gamma} = f_{\gamma \alpha}(y_{\alpha}), x_{\gamma}^{\prime} = f_{\gamma \beta}(x_{\beta}^{\prime}), y_{\gamma}^{\prime} = f_{\gamma \beta}(y_{\beta}^{\prime})\) . By the definition of direct limit, there exists \(\gamma\) in I with \(\gamma \geqslant \alpha, \gamma \geqslant \beta, x_{\gamma} = x_{\gamma}^{\prime}, y_{\gamma} = y_{\gamma}^{\prime}\) . Then

\[
\begin{array}{r l} f _ {\alpha} (x _ {\alpha} y _ {\alpha}) & = f _ {\gamma} (f _ {\gamma \alpha} (x _ {\alpha} y _ {\alpha})) = f _ {\gamma} (x _ {\gamma} y _ {\gamma}) = f _ {\gamma} (x _ {\gamma} ^ {\prime} y _ {\gamma} ^ {\prime}) = f _ {\gamma} (f _ {\gamma \beta} (x _ {\beta} ^ {\prime} y _ {\beta} ^ {\prime})) \\ & = f _ {\beta} (x _ {\beta} ^ {\prime} y _ {\beta} ^ {\prime}). \end{array}
\]

Suppose the magmas \(E_{\alpha}\) are associative. Let x, y, z be in E. There exist \(\alpha \in I\) and elements \(x_{\alpha}, y_{\alpha}, z_{\alpha}\) in \(E_{\alpha}\) such that

\[
x = f _ {\alpha} (x _ {\alpha}), \quad y = f _ {\alpha} (y _ {\alpha}), \quad z = f _ {\alpha} (z _ {\alpha}).
\]

Then \(xy = f_{\alpha}(x_{\alpha}y_{\alpha})\) , whence \((xy)z = f_{\alpha}((x_{\alpha}y_{\alpha})z_{\alpha})\) ; similarly

\[
x (y z) = f _ {\alpha} \left(x _ {\alpha} \left(y _ {\alpha} z _ {\alpha}\right)\right),
\]

whence \((xy)z = x(yz)\) for \((x_{\alpha}y_{\alpha})z_{\alpha} = x_{\alpha}(y_{\alpha}z_{\alpha})\) . The case of commutative magmas is treated analogously.

Suppose finally that each magma \(E_{\alpha}\) has an identity element \(e_{\alpha}\) and that \(f_{\beta\alpha}(e_{\alpha}) = e_{\beta}\) for \(\alpha \leqslant \beta\) . For \(\alpha, \beta\) in I, there exists \(\gamma\) in I with \(\gamma \geqslant \alpha\) and \(\gamma \geqslant \beta\) , whence

\[
f _ {\alpha} (e _ {\alpha}) = f _ {\gamma} \left(f _ {\gamma \alpha} (e _ {\alpha})\right) = f _ {\gamma} (e _ {\gamma}) = f _ {\gamma} \left(f _ {\gamma \beta} (e _ {\beta})\right) = f _ {\beta} (e _ {\beta})
\]

and there thus exists an element \(e\) in \(\mathbf{E}\) such that \(f_{\alpha}(e_{\alpha}) = e\) for all \(\alpha \in \mathbf{I}\) . Let \(x\) be in \(\mathbf{E}\) ; choose \(\alpha \in \mathbf{I}\) and \(x_{\alpha} \in \mathrm{E}_{\alpha}\) such that \(x = f_{\alpha}(x_{\alpha})\) . Then

\[
e x = f _ {\alpha} (e _ {\alpha}) \cdot f _ {\alpha} (x _ {\alpha}) = f _ {\alpha} (e _ {\alpha}. x _ {\alpha}) = f _ {\alpha} (x _ {\alpha}) = x
\]

and similarly \(x.e = x\) , hence \(e\) is the identity element of \(E\) .

The magma E is called the direct limit of the magmas \(E_{\alpha}\) .

PROPOSITION 2. Let \((\mathrm{E}_{\alpha}, f_{\beta \alpha})\) be a direct system of magmas and let E be its direct limits \(f_{\alpha} \colon \mathrm{E}_{\alpha} \to \mathrm{E}\) the canonical homomorphisms. Suppose a magma F and homomorphisms \(u_{\alpha} \colon \mathrm{E}_{\alpha} \to \mathrm{F}\) are given such that \(u_{\alpha} = u_{\beta} \circ f_{\beta \alpha}\) for \(\alpha \leqslant \beta\) . There exists one and only one homomorphism \(u \colon \mathrm{E} \to \mathrm{F}\) such that \(u_{\alpha} = u \circ f_{\alpha}\) for all \(\alpha \in \mathbf{I}\) . If the magmas \(\mathrm{E}_{\alpha}\) and F and the homomorphisms \(f_{\beta \alpha}\) and \(u_{\alpha}\) are unital, the homomorphism \(u\) is unital.

We know (Set Theory, III, § 7, no. 6, Proposition 6) that there exists one and only one mapping \(u: \mathrm{E} \to \mathrm{F}\) such that \(u_{\alpha} = u \circ f_{\alpha}\) for all \(\alpha \in \mathbf{I}\) . We verify that \(u\) is a homomorphism: let \(x, y\) be in \(\mathrm{E}, \alpha\) in \(\mathrm{I}\) and \(x_{\alpha}, y_{\alpha}\) in \(\mathrm{E}_{\alpha}\) such that \(x = f_{\alpha}(x_{\alpha})\) and \(y = f_{\alpha}(y_{\alpha})\) . Then \(xy = f_{\alpha}(x_{\alpha}y_{\alpha})\) , whence

\[
\begin{array}{r l} u (x y) & = u (f _ {\alpha} (x _ {\alpha} y _ {\alpha})) = u _ {\alpha} (x _ {\alpha} y _ {\alpha}) = u _ {\alpha} (x _ {\alpha}) u _ {\alpha} (y _ {\alpha}) \\ & = u (f _ {\alpha} (u _ {\alpha})) u (f _ {\alpha} (y _ {\alpha})) = u (x) u (y). \end{array}
\]

We consider now the unital case and let \(e_{\alpha}\) denote the identity element of \(E_{\alpha}\) , \(e\) that of \(E\) and \(e'\) that of \(F\) . Choose \(\alpha \in I\) , then \(e = f_{\alpha}(e_{\alpha})\) , whence

\[
u (e) = u \left(f _ {\alpha} \left(e _ {\alpha}\right)\right) = u _ {\alpha} \left(e _ {\alpha}\right) = e ^ {\prime}
\]

for \(u_{\alpha}\) is unital. Hence \(u\) is unital.

By analogy with the notion of direct system of magmas, that of a direct system of monoids or groups can be formulated. Proposition 1 shows that the magma E which is the limit of a direct system of monoids \((\mathrm{E}_{\alpha}, f_{\beta \alpha})_{\alpha, \beta \in I}\) is a monoid. We show that E is a group if the \(\mathrm{E}_{\alpha}\) are groups: let \(x \in \mathrm{E}\) , \(\alpha \in I\) and \(x_{\alpha} \in \mathrm{E}_{\alpha}\) be such that \(x = f_{\alpha}(x_{\alpha})\) ; the element \(y = f_{\alpha}(x_{\alpha}^{-1})\) of E is the inverse of \(x\) (§ 2, no. 3). The universal property of Proposition 2 goes over immediately in the case of a direct system of monoids or groups.

The reader is left to define a direct system of rings. Let \((\mathbf{A}_{\alpha}, f_{\beta \alpha})\) be such a direct system; let \(\mathbf{A} = \lim_{\longrightarrow} \mathbf{A}_{\alpha}\) and \(f_{\alpha}: \mathbf{A}_{\alpha} \to \mathbf{A}\) the canonical homomorphisms. There exists (Proposition 2) on \(\mathbf{A}\) one and only one addition and multiplication characterized by \(x + y = f_{\alpha}(x_{\alpha} + y_{\alpha})\) , \(xy = f_{\alpha}(x_{\alpha}y_{\alpha})\) for \(\alpha \in \mathbf{I}\) , \(x_{\alpha}, y_{\alpha}\) in \(\mathbf{A}_{\alpha}\) and \(x = f_{\alpha}(x_{\alpha})\) , \(y = f_{\alpha}(y_{\alpha})\) . Under addition \(\mathbf{A}\) is a commutative group and multiplication is associative and unital. Finally, for \(x, y, z\) in \(\mathbf{A}\) , choose \(\alpha\) in \(\mathbf{I}\) and \(x_{\alpha}, y_{\alpha}\) and \(z_{\alpha}\) in \(\mathbf{A}_{\alpha}\) with

\[
x = f _ {\alpha} (x _ {\alpha}), \quad y = f _ {\alpha} (y _ {\alpha}) \quad \text { and } \quad z = f _ {\alpha} (z _ {\alpha}).
\]

Then

\[
\begin{array}{r l} (x + y) \cdot z & = f _ {\alpha} (x _ {\alpha} + y _ {\alpha}) f _ {\alpha} (z _ {\alpha}) = f _ {\alpha} ((x _ {\alpha} + y _ {\alpha}) z _ {\alpha}) \\ & = f _ {\alpha} (x _ {\alpha} z _ {\alpha} + y _ {\alpha} z _ {\alpha}) = f _ {\alpha} (x _ {\alpha} z _ {\alpha}) + f _ {\alpha} (y _ {\alpha} z _ {\alpha}) = x z + y z \end{array}
\]

and the relation \(x(y + z) = xy + xz\) is proved analogously. In other words, A has the structure of a ring, characterized by the fact that \(f_{\alpha}\) is a ring homomorphism for all \(\alpha \in I\) .

The ring A is called the direct limit of the rings \(A_{\alpha}\) . Proposition 2 extends immediately to the case of rings.

PROPOSITION 3. (a) If the \(\mathbf{A}_{\alpha}\) are non-zero, \(\mathbf{A}\) is non-zero.

\begin{enumerate}
\def\labelenumi{(\alph{enumi})}
\setcounter{enumi}{1}
\item
  If the \(\mathbf{A}_{\alpha}\) are integral domains, \(\mathbf{A}\) is an integral domain.
\item
  If the \(\mathbf{A}_{\alpha}\) are fields, \(\mathbf{A}\) is a field.
\end{enumerate}

Let \(0_{\alpha}, 1_{\alpha}\) be the zero and unit of \(A_{\alpha}\) and 0, 1 the zero and unit of A. There exists \(\alpha \in I\) such that \(f_{\alpha}(0_{\alpha}) = 0, f_{\alpha}(1_{\alpha}) = 1\) . If 0 = 1, there exists \(\beta \geqslant \alpha\) such that \(f_{\beta \alpha}(0_{\alpha}) = f_{\beta \alpha}(1_{\alpha})\) , that is \(0_{\beta} = 1_{\beta}\) . This proves (a).

Suppose that \(A_{\alpha}\) are integral domains. Then A is commutative and non-zero by (a). Let x, y be elements of A such that xy = 0. There exists \(\alpha \in I\) and \(x_{\alpha}, y_{\alpha} \in A_{\alpha}\) such that \(x = f_{\alpha}(x_{\alpha}), y = f_{\alpha}(y_{\alpha})\) . Then \(f_{\alpha}(x_{\alpha}y_{\alpha}) = xy = 0 = f_{\alpha}(0_{\alpha})\) . Hence there exists \(\beta \geqslant \alpha\) such that \(f_{\beta\alpha}(x_{\alpha}y_{\alpha}) = f_{\beta\alpha}(0_{\alpha})\) . As \(A_{\beta}\) is an integral domain, it follows that \(f_{\beta\alpha}(x_{\alpha}) = 0_{\beta}\) or \(f_{\beta\alpha}(y_{\alpha}) = 0_{\beta}\) , hence x = 0 or y = 0. This proves (b).

Suppose that the \(A_{\alpha}\) are fields. Then \(A \neq \{0\}\) by (a). Let x be a non-zero element of A. There exist \(\alpha \in I\) and \(x_{\alpha} \in A_{\alpha}\) such that \(x = f_{\alpha}(x_{\alpha})\) . Then \(x_{\alpha} \neq 0\) and \(f_{\alpha}(x_{\alpha}^{-1})\) is the inverse of x in A. This proves (c).

Let \(\mathfrak{E} = (\mathrm{E}_{\alpha}, f_{\beta \alpha})\) and \(\mathfrak{E}' = (\mathrm{E}_{\alpha}', f_{\beta \alpha}')\) be two direct systems of magmas (resp. monoids, groups, rings). A homomorphism of \(\mathfrak{E}\) into \(\mathfrak{E}'\) is a direct system \((u_{\alpha})_{\alpha \in I}\) of mappings \(u_{\alpha} \colon E_{\alpha} \to E_{\alpha}'\) such that each \(u_{\alpha}\) is a homomorphism. Under these conditions, the mapping \(u = \lim u_{\alpha}\) from \(E = \lim E_{\alpha}\) to \(E' = \lim E_{\alpha}'\) is a homomorphism (cf.~Set Theory, III, § 7, no. 6).

